\section{Related Work}
\label{sec:related_work}

\begin{figure*}[t]
    \centering
    \includegraphics[width=1\linewidth]{figures/overview.pdf}
    \caption{\textbf{Overview of RT-DETRv4}. We leverage a Vision Foundation Model (VFM) to extract high-quality semantic representations, which are aligned with the deepest feature map ($F_5$) from the AIFI module via a Feature Projector in the Deep Semantic Injector (DSI). To ensure faster and more stable convergence, a Gradient-guided Adaptive Modulation (GAM) dynamically adjusts the DSI loss during training. The proposed framework operates only during the training phase (highlighted by dashed arrows and blue blocks) of the real-time detector and keeps the original architecture unchanged during inference and deployment, introducing no additional overhead while improving accuracy.}
    \label{fig:overview}
\vspace{-5mm}
\end{figure*}

\subsection{Real-time Object Detection}
The evolution of real-time object detection has long been driven by the You Only Look Once (YOLO) family~\cite{redmon2016you}, which popularized the single-stage paradigm through an efficient and unified detection pipeline. Over the past few years, this lineage has undergone rapid iteration, introducing continuous refinements in backbone design, label assignment, and optimization strategy~\cite{redmon2017yolo9000, redmon2018yolov3, bochkovskiy2020yolov4, yolov5v7.0, li2022yolov6, wang2023yolov7, yolov8, wang2024yolov9}.
Recent generations have expanded the design space even further: YOLOv10~\cite{wang2024yolov10} eliminated NMS that the YOLO series has long relied on, YOLO11~\cite{yolo11} improved the architectural hierarchy and neck connectivity, YOLOv12~\cite{tian2025yolov12} incorporated attention mechanisms for better contextual reasoning, and YOLOv13~\cite{lei2025yolov13} explored hypergraph representations to capture higher-order feature dependencies.
These advances have pushed the performance–efficiency frontier of convolutional and hybrid architectures, gradually narrowing the gap between real-time and high-accuracy detectors.

In parallel, another line of research has evolved around the DEtection TRansformer (DETR)~\cite{carion2020end}, which redefined object detection as a set prediction problem and eliminated hand-crafted components such as anchor design and NMS. This transformer-based paradigm inspired numerous variants, including Deformable DETR~\cite{zhu2020deformable}, Conditional DETR~\cite{meng2021conditional}, and DAB-DETR~\cite{liu2022dab}, which focus on improving convergence and localization accuracy. Later works such as DN-DETR~\cite{li2022dn}, DINO~\cite{zhang2022dino}, and Group-DETR~\cite{chen2023group} introduced denoising objectives and group-wise supervision to further enhance training stability and representational quality.

Building on this foundation, RT-DETR~\cite{zhao2024detrs} established the first real-time end-to-end transformer detector that achieved parity with, and in some cases surpassed, contemporary YOLO models. Subsequent works have continued to improve its training efficiency and representation learning without incurring inference overhead. For instance, RT-DETRv2~\cite{lv2024rt} and RT-DETRv3~\cite{wang2025rt} incorporated auxiliary supervision for enhanced gradient flow, D-FINE~\cite{peng2024d} employed self-distillation to refine semantic representation, and DEIM~\cite{huang2025deim} introduced dense matching for more precise feature alignment.
Collectively, these developments illustrate a clear trend: as architectural efficiency saturates, training supervision and semantic representation become the primary levers for further progress. Our work builds on this insight by strengthening the core representation via deep semantic transfer, achieving higher accuracy at no additional deployment or inference cost.

\subsection{Vision Foundation Models}
Vision Foundation Models (VFMs) have become a dominant paradigm for learning general-purpose vision representations from large-scale image corpora with minimal or no human supervision. Early progress stemmed from self- and weakly-supervised learning methods, which enabled models to capture high-level semantics from unlabeled or loosely labeled data. Representative approaches include contrastive learning frameworks such as SimCLR~\cite{chen2020simple} and MoCo~\cite{he2020momentum}, which learn discriminative features by enforcing consistency across augmented views of the same image while contrasting them with others. CLIP~\cite{radford2021learning} further extended this idea to large-scale image–text contrastive training, aligning visual and linguistic embeddings and demonstrating strong zero-shot transferability across diverse tasks.

Inspired by masked language modeling in NLP, Masked Image Modeling (MIM) approaches were introduced to reconstruct masked image regions, thereby learning context-aware and holistic representations. Notable methods include MAE~\cite{he2022masked} and BEiT~\cite{baobeit}. Building on these advances, the DINO family~\cite{caron2021emerging, oquab2023dinov2, simeoni2025dinov3} integrates contrastive, reconstruction-based, and self-distillation objectives to produce highly semantic and transferable features. In particular, DINOv3~\cite{simeoni2025dinov3} demonstrates the scalability and efficacy of large-scale self-supervised learning, achieving rich and robust representations without human annotations.